\begin{table}[H]
  \centering
  \small
  \caption{Nhãn policy hoà điểm trong hai tập dữ liệu chuyên gia}
  \label{tab:cs-dl-hoadiem}
  \begin{tabular}{lrr}
    \toprule
    Đại lượng & expert\_v1 & expert\_v2 \\
    \midrule
    Số mẫu & 27.720 & 297.098 \\
    Mẫu có nhãn hoà điểm & 12.601 & 48.043 \\
    Mẫu ở thế đã thua & 902 & 35.045 \\
    Thế đã thua có nhãn hoà & 100,0\% & 100,0\% \\
    Nước nhãn cao nhất ở hàng 1: toàn tập & 17,0\% & 14,7\% \\
    \quad trong mẫu hoà & 31,0\% & 73,8\% \\
    \quad trong mẫu không hoà & 5,4\% & 3,3\% \\
    \quad trong thế đã thua & 82,6\% & 91,9\% \\
    Nhãn hàng 1 đến từ mẫu hoà & 82,6\% & 81,3\% \\
    Ván không còn mẫu sau loại trùng & 17.242 & 1.178 \\
    Ván còn không quá 2 mẫu & 19.009 & 3.643 \\
    \bottomrule
  \end{tabular}
  \par\smallskip\parbox{.95\linewidth}{\footnotesize Hoà điểm: nhãn đặt xác suất bằng nhau lên từ hai nước trở lên. Thế đã thua: đối thủ có từ hai ô thắng, bên đi không có ô thắng.}
\end{table}
